% --- Archon physics formalization source begin ---
% archon:physics
% archon:covers PhyXMiniProblems/problem_phyx_mini_0312.lean
% archon:source-report reports/phyx_mini/problem_phyx_mini_0312.source.json

\chapter{Physics problem phyx\_mini\_0312}
\label{ch:PhyXMiniProblems_problem_phyx_mini_0312}

\paragraph{Problem source.}
\#\# Physical scenario

Figure shows two isotropic point sources of sound, \$S\_1\$ and \$S\_2\$. The sources emit waves in phase at wavelength \$0.50\$ m; they are separated by \$D = 1.75\$ m.

\#\# Question

If we move a sound detector along a large circle centered at the midpoint between the sources, at how many points do waves arrive at the detector exactly out of phase?

\#\# Answer choices

- A: A: 11

- B: B: 12

- C: C: 13

- D: D: 14

\#\# Figure caption (auxiliary; use the image as primary evidence)

The image depicts a horizontal line with two red dots positioned on it, labeled as \textbackslash{}(S\_1\textbackslash{}) and \textbackslash{}(S\_2\textbackslash{}). The dots are positioned at different points along the line, and arrows point away from \textbackslash{}(S\_1\textbackslash{}) toward \textbackslash{}(S\_2\textbackslash{}), indicating direction. The arrows are accompanied by the letter \textbackslash{}(D\textbackslash{}), suggesting a distance or relationship between \textbackslash{}(S\_1\textbackslash{}) and \textbackslash{}(S\_2\textbackslash{}). The arrows reinforce the concept that \textbackslash{}(D\textbackslash{}) represents a measurement or parameter between the two points.

\#\# Dataset metadata

Resonance and Harmonics; Physical Model Grounding Reasoning, Spatial Relation Reasoning

\paragraph{Recorded answer/context.}
D: D: 14

\paragraph{Figure/image path.}
/root/proposal\_for\_physic/science-mango/phyx\_mini\_run/phyx\_data/test\_image/312.png

\paragraph{Formalization target.}
create a compiling Lean file with sorry bodies at `PhyXMiniProblems/problem\_phyx\_mini\_0312.lean`. The Lean declarations must preserve the physical quantities, dimensions or dimensional roles, figure labels, governing-law hypotheses, and final relation expressed by this problem.
Use Mathlib/Physlib names found through LeanExplore where available. If a physics API is missing, introduce faithful local abstractions rather than scalar placeholder aliases.

\begin{theorem}[Physics formalization target]
\label{thm:physics:phyx_mini_0312:target}
\lean{PhyXMiniProblems.ProblemPhyXMini0312.problem_phyx_mini_0312}
\uses{def:physics:phyx-mini-0312:phyxminiproblems-problemphyxmini0312-twopointsourceinterferencesetup, def:physics:phyx-mini-0312:phyxminiproblems-problemphyxmini0312-matchesproblemreadouts, def:physics:phyx-mini-0312:phyxminiproblems-problemphyxmini0312-matchescoherentisotropicsourcescenario, def:physics:phyx-mini-0312:phyxminiproblems-problemphyxmini0312-matchessuppliedfigure, def:physics:phyx-mini-0312:phyxminiproblems-problemphyxmini0312-matchessourcegeometry, def:physics:phyx-mini-0312:phyxminiproblems-problemphyxmini0312-hasphysicallargecircleparameters, def:physics:phyx-mini-0312:phyxminiproblems-problemphyxmini0312-satisfiestwopointsourceinterferencelaw, def:physics:phyx-mini-0312:phyxminiproblems-problemphyxmini0312-outofphasedetectorpoints, def:physics:phyx-mini-0312:phyxminiproblems-problemphyxmini0312-isuniquematchinganswer}
The assigned autoformalize agent should translate this physics problem into Lean declarations in the covered file, with theorem and lemma proof bodies written as `by sorry`.
\end{theorem}
\begin{proof}
This is an autoformalization task, not a proof task. Produce statements that can later be proved by the physics prover without weakening the physical model.
\end{proof}

% --- Archon declaration topology begin ---
\section{Declaration topology}
\begin{definition}[Length Quantity]
  \label{def:physics:phyx-mini-0312:phyxminiproblems-problemphyxmini0312-lengthquantity}
  \lean{PhyXMiniProblems.ProblemPhyXMini0312.LengthQuantity}
  A nonnegative, unit-independent physical length
\end{definition}
\begin{proof}
  This is the declared model definition of Length Quantity.
\end{proof}
\begin{definition}[length Readout]
  \label{def:physics:phyx-mini-0312:phyxminiproblems-problemphyxmini0312-lengthreadout}
  \lean{PhyXMiniProblems.ProblemPhyXMini0312.lengthReadout}
  \uses{def:physics:phyx-mini-0312:phyxminiproblems-problemphyxmini0312-lengthquantity}
  Read a physical length as a real number in a selected length unit
\end{definition}
\begin{proof}
  This is the declared model definition of length Readout.
\end{proof}
\begin{definition}[length In Meters]
  \label{def:physics:phyx-mini-0312:phyxminiproblems-problemphyxmini0312-lengthinmeters}
  \lean{PhyXMiniProblems.ProblemPhyXMini0312.lengthInMeters}
  \uses{def:physics:phyx-mini-0312:phyxminiproblems-problemphyxmini0312-lengthquantity, def:physics:phyx-mini-0312:phyxminiproblems-problemphyxmini0312-lengthreadout}
  Metre readout of a physical length
\end{definition}
\begin{proof}
  This is the declared model definition of length In Meters.
\end{proof}
\begin{definition}[Source Label]
  \label{def:physics:phyx-mini-0312:phyxminiproblems-problemphyxmini0312-sourcelabel}
  \lean{PhyXMiniProblems.ProblemPhyXMini0312.SourceLabel}
  The two source labels printed in the supplied figure
\end{definition}
\begin{proof}
  This is the declared model definition of Source Label.
\end{proof}
\begin{definition}[Acoustic Source Kind]
  \label{def:physics:phyx-mini-0312:phyxminiproblems-problemphyxmini0312-acousticsourcekind}
  \lean{PhyXMiniProblems.ProblemPhyXMini0312.AcousticSourceKind}
  The idealized acoustic-source kind stipulated by the problem
\end{definition}
\begin{proof}
  This is the declared model definition of Acoustic Source Kind.
\end{proof}
\begin{definition}[Figure Label]
  \label{def:physics:phyx-mini-0312:phyxminiproblems-problemphyxmini0312-figurelabel}
  \lean{PhyXMiniProblems.ProblemPhyXMini0312.FigureLabel}
  Literal symbolic labels visible in the supplied raster
\end{definition}
\begin{proof}
  This is the declared model definition of Figure Label.
\end{proof}
\begin{definition}[Source Separation Figure]
  \label{def:physics:phyx-mini-0312:phyxminiproblems-problemphyxmini0312-sourceseparationfigure}
  \lean{PhyXMiniProblems.ProblemPhyXMini0312.SourceSeparationFigure}
  \uses{def:physics:phyx-mini-0312:phyxminiproblems-problemphyxmini0312-sourcelabel, def:physics:phyx-mini-0312:phyxminiproblems-problemphyxmini0312-figurelabel}
  Qualitative evidence read from the image: two red source dots and their labels lie on a common horizontal baseline, with vertical projections down to a double-headed separation indicator labelled D. The raster itself contains no numerical length
\end{definition}
\begin{proof}
  This is the declared model definition of Source Separation Figure.
\end{proof}
\begin{definition}[Two Point Source Interference Setup]
  \label{def:physics:phyx-mini-0312:phyxminiproblems-problemphyxmini0312-twopointsourceinterferencesetup}
  \lean{PhyXMiniProblems.ProblemPhyXMini0312.TwoPointSourceInterferenceSetup}
  \uses{def:physics:phyx-mini-0312:phyxminiproblems-problemphyxmini0312-lengthquantity, def:physics:phyx-mini-0312:phyxminiproblems-problemphyxmini0312-sourcelabel, def:physics:phyx-mini-0312:phyxminiproblems-problemphyxmini0312-acousticsourcekind, def:physics:phyx-mini-0312:phyxminiproblems-problemphyxmini0312-sourceseparationfigure}
  The independent quantities, positions, emission phases, and arrival-phase observable of the experiment. The arrival observable is not defined from the requested count; it is constrained only by the governing interference law
\end{definition}
\begin{proof}
  This is the declared model definition of Two Point Source Interference Setup.
\end{proof}
\begin{definition}[Matches Problem Readouts]
  \label{def:physics:phyx-mini-0312:phyxminiproblems-problemphyxmini0312-matchesproblemreadouts}
  \lean{PhyXMiniProblems.ProblemPhyXMini0312.MatchesProblemReadouts}
  \uses{def:physics:phyx-mini-0312:phyxminiproblems-problemphyxmini0312-lengthinmeters, def:physics:phyx-mini-0312:phyxminiproblems-problemphyxmini0312-twopointsourceinterferencesetup}
  Numerical physical readouts stated in the problem prose
\end{definition}
\begin{proof}
  This is the declared model definition of Matches Problem Readouts.
\end{proof}
\begin{definition}[Matches Coherent Isotropic Source Scenario]
  \label{def:physics:phyx-mini-0312:phyxminiproblems-problemphyxmini0312-matchescoherentisotropicsourcescenario}
  \lean{PhyXMiniProblems.ProblemPhyXMini0312.MatchesCoherentIsotropicSourceScenario}
  \uses{def:physics:phyx-mini-0312:phyxminiproblems-problemphyxmini0312-sourcelabel, def:physics:phyx-mini-0312:phyxminiproblems-problemphyxmini0312-twopointsourceinterferencesetup}
  Both labelled objects are isotropic point sources of sound and have equal initial emission phase
\end{definition}
\begin{proof}
  This is the declared model definition of Matches Coherent Isotropic Source Scenario.
\end{proof}
\begin{definition}[Matches Supplied Figure]
  \label{def:physics:phyx-mini-0312:phyxminiproblems-problemphyxmini0312-matchessuppliedfigure}
  \lean{PhyXMiniProblems.ProblemPhyXMini0312.MatchesSuppliedFigure}
  \uses{def:physics:phyx-mini-0312:phyxminiproblems-problemphyxmini0312-sourcelabel, def:physics:phyx-mini-0312:phyxminiproblems-problemphyxmini0312-figurelabel, def:physics:phyx-mini-0312:phyxminiproblems-problemphyxmini0312-twopointsourceinterferencesetup}
  Labels, colors, and incidences transcribed from the primary image
\end{definition}
\begin{proof}
  This is the declared model definition of Matches Supplied Figure.
\end{proof}
\begin{definition}[Matches Source Geometry]
  \label{def:physics:phyx-mini-0312:phyxminiproblems-problemphyxmini0312-matchessourcegeometry}
  \lean{PhyXMiniProblems.ProblemPhyXMini0312.MatchesSourceGeometry}
  \uses{def:physics:phyx-mini-0312:phyxminiproblems-problemphyxmini0312-lengthreadout, def:physics:phyx-mini-0312:phyxminiproblems-problemphyxmini0312-twopointsourceinterferencesetup}
  Planar geometry corresponding to the figure and the phrase "midpoint between the sources." Coordinate 0 is horizontal and coordinate 1 is vertical. The metric and dimensionful separation use the same selected length unit
\end{definition}
\begin{proof}
  This is the declared model definition of Matches Source Geometry.
\end{proof}
\begin{definition}[Has Physical Large Circle Parameters]
  \label{def:physics:phyx-mini-0312:phyxminiproblems-problemphyxmini0312-hasphysicallargecircleparameters}
  \lean{PhyXMiniProblems.ProblemPhyXMini0312.HasPhysicalLargeCircleParameters}
  \uses{def:physics:phyx-mini-0312:phyxminiproblems-problemphyxmini0312-lengthreadout, def:physics:phyx-mini-0312:phyxminiproblems-problemphyxmini0312-twopointsourceinterferencesetup}
  Strictly positive lengths and a precise form of "large circle": the radius is greater than half the source separation, so both sources lie strictly inside the circle once the midpoint geometry is imposed. No requested point count is included here
\end{definition}
\begin{proof}
  This is the declared model definition of Has Physical Large Circle Parameters.
\end{proof}
\begin{definition}[detector Circle]
  \label{def:physics:phyx-mini-0312:phyxminiproblems-problemphyxmini0312-detectorcircle}
  \lean{PhyXMiniProblems.ProblemPhyXMini0312.detectorCircle}
  \uses{def:physics:phyx-mini-0312:phyxminiproblems-problemphyxmini0312-lengthreadout, def:physics:phyx-mini-0312:phyxminiproblems-problemphyxmini0312-twopointsourceinterferencesetup}
  The large circular path along which the detector is moved
\end{definition}
\begin{proof}
  This is the declared model definition of detector Circle.
\end{proof}
\begin{definition}[path Difference]
  \label{def:physics:phyx-mini-0312:phyxminiproblems-problemphyxmini0312-pathdifference}
  \lean{PhyXMiniProblems.ProblemPhyXMini0312.pathDifference}
  \uses{def:physics:phyx-mini-0312:phyxminiproblems-problemphyxmini0312-twopointsourceinterferencesetup}
  Absolute difference between the two source-to-detector path lengths
\end{definition}
\begin{proof}
  This is the declared model definition of path Difference.
\end{proof}
\begin{definition}[Is Odd Half Wavelength Path Difference]
  \label{def:physics:phyx-mini-0312:phyxminiproblems-problemphyxmini0312-isoddhalfwavelengthpathdifference}
  \lean{PhyXMiniProblems.ProblemPhyXMini0312.IsOddHalfWavelengthPathDifference}
  \uses{def:physics:phyx-mini-0312:phyxminiproblems-problemphyxmini0312-lengthreadout, def:physics:phyx-mini-0312:phyxminiproblems-problemphyxmini0312-twopointsourceinterferencesetup, def:physics:phyx-mini-0312:phyxminiproblems-problemphyxmini0312-pathdifference}
  For coherent in-phase sources, destructive arrival occurs at an odd half-integral multiple (2 n + 1) λ / 2 of the common wavelength. Natural orders index the nonnegative absolute path differences
\end{definition}
\begin{proof}
  This is the declared model definition of Is Odd Half Wavelength Path Difference.
\end{proof}
\begin{definition}[Satisfies Two Point Source Interference Law]
  \label{def:physics:phyx-mini-0312:phyxminiproblems-problemphyxmini0312-satisfiestwopointsourceinterferencelaw}
  \lean{PhyXMiniProblems.ProblemPhyXMini0312.SatisfiesTwoPointSourceInterferenceLaw}
  \uses{def:physics:phyx-mini-0312:phyxminiproblems-problemphyxmini0312-twopointsourceinterferencesetup, def:physics:phyx-mini-0312:phyxminiproblems-problemphyxmini0312-isoddhalfwavelengthpathdifference}
  The governing two-point-source phase law. It relates the independent arrival observable to path difference and wavelength at every planar detector point; it contains neither a cardinality nor an answer label
\end{definition}
\begin{proof}
  This is the declared model definition of Satisfies Two Point Source Interference Law.
\end{proof}
\begin{definition}[out Of Phase Detector Points]
  \label{def:physics:phyx-mini-0312:phyxminiproblems-problemphyxmini0312-outofphasedetectorpoints}
  \lean{PhyXMiniProblems.ProblemPhyXMini0312.outOfPhaseDetectorPoints}
  \uses{def:physics:phyx-mini-0312:phyxminiproblems-problemphyxmini0312-twopointsourceinterferencesetup, def:physics:phyx-mini-0312:phyxminiproblems-problemphyxmini0312-detectorcircle}
  Points on the detector circle where the two waves arrive in opposition
\end{definition}
\begin{proof}
  This is the declared model definition of out Of Phase Detector Points.
\end{proof}
\begin{definition}[Answer Choice]
  \label{def:physics:phyx-mini-0312:phyxminiproblems-problemphyxmini0312-answerchoice}
  \lean{PhyXMiniProblems.ProblemPhyXMini0312.AnswerChoice}
  Labels of the four printed answer choices
\end{definition}
\begin{proof}
  This is the declared model definition of Answer Choice.
\end{proof}
\begin{definition}[displayed Point Count]
  \label{def:physics:phyx-mini-0312:phyxminiproblems-problemphyxmini0312-displayedpointcount}
  \lean{PhyXMiniProblems.ProblemPhyXMini0312.displayedPointCount}
  \uses{def:physics:phyx-mini-0312:phyxminiproblems-problemphyxmini0312-answerchoice}
  Detector-point counts printed beside the answer labels
\end{definition}
\begin{proof}
  This is the declared model definition of displayed Point Count.
\end{proof}
\begin{definition}[recorded Dataset Answer]
  \label{def:physics:phyx-mini-0312:phyxminiproblems-problemphyxmini0312-recordeddatasetanswer}
  \lean{PhyXMiniProblems.ProblemPhyXMini0312.recordedDatasetAnswer}
  \uses{def:physics:phyx-mini-0312:phyxminiproblems-problemphyxmini0312-answerchoice}
  The answer label recorded in the source dataset
\end{definition}
\begin{proof}
  This is the declared model definition of recorded Dataset Answer.
\end{proof}
\begin{definition}[Matches Displayed Point Count]
  \label{def:physics:phyx-mini-0312:phyxminiproblems-problemphyxmini0312-matchesdisplayedpointcount}
  \lean{PhyXMiniProblems.ProblemPhyXMini0312.MatchesDisplayedPointCount}
  \uses{def:physics:phyx-mini-0312:phyxminiproblems-problemphyxmini0312-twopointsourceinterferencesetup, def:physics:phyx-mini-0312:phyxminiproblems-problemphyxmini0312-outofphasedetectorpoints, def:physics:phyx-mini-0312:phyxminiproblems-problemphyxmini0312-answerchoice, def:physics:phyx-mini-0312:phyxminiproblems-problemphyxmini0312-displayedpointcount}
  A choice agrees with the cardinality derived from the physical model
\end{definition}
\begin{proof}
  This is the declared model definition of Matches Displayed Point Count.
\end{proof}
\begin{definition}[Is Unique Matching Answer]
  \label{def:physics:phyx-mini-0312:phyxminiproblems-problemphyxmini0312-isuniquematchinganswer}
  \lean{PhyXMiniProblems.ProblemPhyXMini0312.IsUniqueMatchingAnswer}
  \uses{def:physics:phyx-mini-0312:phyxminiproblems-problemphyxmini0312-twopointsourceinterferencesetup, def:physics:phyx-mini-0312:phyxminiproblems-problemphyxmini0312-answerchoice, def:physics:phyx-mini-0312:phyxminiproblems-problemphyxmini0312-matchesdisplayedpointcount}
  Exactly one printed choice agrees with the physical detector count
\end{definition}
\begin{proof}
  This is the declared model definition of Is Unique Matching Answer.
\end{proof}
\begin{lemma}[out Of Phase Detector Point Count exact]
  \label{lem:physics:phyx-mini-0312:phyxminiproblems-problemphyxmini0312-outofphasedetectorpointcount-exact}
  \lean{PhyXMiniProblems.ProblemPhyXMini0312.outOfPhaseDetectorPointCount_exact}
  \uses{def:physics:phyx-mini-0312:phyxminiproblems-problemphyxmini0312-twopointsourceinterferencesetup, def:physics:phyx-mini-0312:phyxminiproblems-problemphyxmini0312-matchesproblemreadouts, def:physics:phyx-mini-0312:phyxminiproblems-problemphyxmini0312-matchescoherentisotropicsourcescenario, def:physics:phyx-mini-0312:phyxminiproblems-problemphyxmini0312-matchessourcegeometry, def:physics:phyx-mini-0312:phyxminiproblems-problemphyxmini0312-hasphysicallargecircleparameters, def:physics:phyx-mini-0312:phyxminiproblems-problemphyxmini0312-satisfiestwopointsourceinterferencelaw, def:physics:phyx-mini-0312:phyxminiproblems-problemphyxmini0312-outofphasedetectorpoints}
  The geometry, readouts, large-circle condition, and odd-half-wavelength phase law give exactly fourteen out-of-phase points. The target count occurs only in this conclusion, not in the setup or governing-law predicates
\end{lemma}
\begin{proof}
  The conclusion follows from the governing relations and figure readouts named in the statement of out Of Phase Detector Point Count exact.
\end{proof}
% --- Archon declaration topology end ---
% --- Archon physics formalization source end ---
